\documentclass[12pt]{article}
\usepackage[utf8]{inputenc}
\usepackage[spanish]{babel}
\usepackage[a4paper, margin=2cm]{geometry}
\usepackage{tikz}
\usetikzlibrary{babel}
\usepackage[most]{tcolorbox}
\usepackage{amsmath}

\renewcommand{\familydefault}{\sfdefault}

\definecolor{medblue}{RGB}{43, 108, 176}
\definecolor{medteal}{RGB}{49, 151, 149}
\definecolor{medlight}{RGB}{235, 248, 255}
\definecolor{meddark}{RGB}{26, 32, 44}
\definecolor{neckpurple}{RGB}{139, 92, 246}

\begin{document}
\pagestyle{empty}

\begin{tcolorbox}[
    enhanced, 
    colback=medlight!30, 
    colframe=medblue, 
    title={\Large \textbf{Ángulo Plano: Terapia Física}}, 
    halign title=center, 
    fonttitle=\bfseries, 
    arc=4mm, 
    boxrule=1.5pt,
    shadow={2mm}{-2mm}{0mm}{black!30!white}
]

    \vspace{0.2cm}
    \begin{tcolorbox}[colback=white, colframe=medteal, arc=2mm, boxrule=1.2pt, title=\textbf{Caso Clínico / Problema}]
        Durante una terapia de cuello, un paciente logra una rotación lateral equivalente a un ángulo de \textbf{$\frac{\pi}{4}$ radianes}. ¿A cuántos grados corresponde esta rotación?
    \end{tcolorbox}

    \vspace{0.4cm}
    
    \begin{center}
    \begin{tikzpicture}[scale=1.5]
        % Guías de la vista superior (cabeza rotando)
        \draw[dashed, meddark!50, line width=1.5pt] (0, -0.5) -- (0, 3.5) node[above, font=\bfseries] {Mirada al frente};
        \draw[dashed, meddark!50, line width=1.5pt] (-2, 0) -- (2, 0) node[right, font=\bfseries] {Hombros};
        
        % Vectores de rotación
        \draw[->, line width=3pt, meddark] (0,0) -- (0, 2.5);
        \draw[->, line width=3pt, neckpurple] (0,0) -- (45:2.5) node[above right, font=\bfseries] {Rotación Máxima};
        
        % Sector del ángulo
        \filldraw[neckpurple, fill opacity=0.2, draw=neckpurple, line width=2pt] (0,0) -- (0,1.5) arc (90:45:1.5) -- cycle;
        
        % Etiqueta del ángulo
        \node[font=\Large\bfseries, text=neckpurple!60!black] at (67.5:2) {$\frac{\pi}{4}$};
        
        % Centro (Cuello)
        \filldraw[meddark] (0,0) circle (4pt) node[below right, font=\bfseries, xshift=2pt, yshift=-2pt] {Cuello (Eje)};
    \end{tikzpicture}
    \end{center}

    \vspace{0.4cm}
    \begin{tcolorbox}[colback=medteal!10, colframe=medteal, arc=2mm, boxrule=1.2pt, title=\textbf{Desarrollo Matemático y Solución}]
        Dado que el ángulo ya se encuentra expresado en función de $\pi$, la conversión a grados sexagesimales resulta muy sencilla multiplicando por \textbf{$\left(\frac{180^\circ}{\pi}\right)$}:
        \begin{align*}
            \theta &= \frac{\pi}{4} \text{ rad} \times \left( \frac{180^\circ}{\pi \text{ rad}} \right) \\[1ex]
            \theta &= \frac{180^\circ \times \pi}{4 \times \pi}
        \end{align*}
        El número $\pi$ se cancela en el numerador y el denominador, por lo que basta dividir 180 entre 4:
        \[ \theta = \frac{180^\circ}{4} = \mathbf{45^\circ} \]
        \textbf{Resultado:} La rotación lateral lograda por el paciente es de \textbf{$45^\circ$}.
    \end{tcolorbox}
    
\end{tcolorbox}

\end{document}